% ==============================================================================
% preamble.tex – gemeinsame Pakete und Befehle für alle Kapitel
% (wird von masterarbeit.tex und Abgabe.tex mit \input geladen)
% ==============================================================================

% --- Kodierung & Sprache -------------------------------------------------------
\usepackage[T1]{fontenc}
\usepackage[utf8]{inputenc}
\usepackage[english,ngerman]{babel}          % letzte Sprache = Hauptsprache
\usepackage[autostyle=true,german=quotes]{csquotes}

% --- Schrift & Typografie ------------------------------------------------------
\usepackage{lmodern}
\usepackage{microtype}
\usepackage{setspace}
\setstretch{1.15}
\setlength{\emergencystretch}{3em}           % vermeidet zu breite Zeilen

% --- Mathematik & Einheiten ----------------------------------------------------
\usepackage{amsmath}
\usepackage{amssymb}
\usepackage{bm}
\usepackage{icomma}                          % Dezimalkomma im Mathemodus
\usepackage[locale=DE]{siunitx}

% --- Tabellen ------------------------------------------------------------------
\usepackage{booktabs}
\usepackage{tabularx}
\usepackage{array}
\usepackage{multirow}
\usepackage{longtable}
\renewcommand{\arraystretch}{1.2}
\newcolumntype{L}[1]{>{\raggedright\arraybackslash}p{#1}}
\newcolumntype{Y}{>{\raggedright\arraybackslash}X}

% --- Abbildungen & Beschriftungen ----------------------------------------------
\usepackage{graphicx}
\graphicspath{{figures/}}
\usepackage{float}
\usepackage[font=small,labelfont=bf]{caption}
\usepackage{subcaption}

% --- Listen, Farben, Hilfsmittel -----------------------------------------------
\usepackage{enumitem}
\usepackage[dvipsnames,table]{xcolor}

% --- Literatur (natbib + bibtex) -----------------------------------------------
\usepackage[numbers,sort&compress]{natbib}

% --- Links & Querverweise (hyperref zuletzt, cleveref danach) -------------------
\usepackage[hidelinks]{hyperref}
\usepackage[ngerman,capitalise,noabbrev]{cleveref}

% --- Eigene Befehle ------------------------------------------------------------
\newcommand{\code}[1]{\texttt{#1}}
\newcommand{\Ediag}{E_{\mathrm{diag}}}
\newcommand{\sel}{\ensuremath{^{\dagger}}}   % gewählte Option in Tabellen
\newcommand{\placeholder}[1]{\textcolor{red}{\textbf{[Noch zu schreiben:} #1\textbf{]}}}
\newcommand{\notebox}[1]{%
  \par\noindent\fcolorbox{gray!60}{gray!8}{\parbox{\dimexpr0.97\linewidth-2\fboxsep-2\fboxrule}{\small #1}}\par}

% --- Angaben zur Arbeit (hier einmal anpassen) ----------------------------------
\newcommand{\Titel}{Dateneffiziente Personalisierung für webcam-basiertes Eye-Tracking mithilfe vortrainierter Bildverarbeitungsmodelle}
\newcommand{\TitelEN}{Data-Efficient Personalization for Webcam-Based Eye Tracking Using Pretrained Vision Models}
\newcommand{\Autor}{Khaled Rafei}
\newcommand{\Matrikel}{930677}
\newcommand{\Studiengang}{Technische Informatik}
\newcommand{\Hochschule}{Berliner Hochschule für Technik}
\newcommand{\Fakultaet}{Fachbereich VI -- Informatik und Medien}
\newcommand{\Arbeitstyp}{Masterarbeit}
\newcommand{\Abschluss}{Master of Engineering (M.\,Eng.)}
% Prüfungskommission laut Zulassung
\newcommand{\Betreuer}{Prof.~Dr.~Kristian Hildebrand}        % Betreuer BHT
\newcommand{\BetreuerExtern}{Prof.~Dr.~Meysam Shamsi}       % Betreuer extern
\newcommand{\Gutachter}{Prof.~Dr.~Robert Strzebkowski}       % Gutachter
\newcommand{\Abgabedatum}{05.\,11.\,2026}

% --- Zweite Sprache (Standard: aus) -------------------------------------------
% \ZweiteSprachefalse: die ganze Arbeit nur auf Deutsch.
% \ZweiteSprachetrue:  zusätzlich englischer Titel auf der Titelseite und ein
%                      englisches Abstract (falls die Prüfungsordnung das verlangt).
\newif\ifZweiteSprache
\ZweiteSprachefalse

% --- BHT-Titelseite (nach dem Muster "Titelblatt_Muster.pdf" der BHT) ------------
% Oben das BHT-Logo (figures/bht/), alles andere schwarz und zentriert.
% Serifenlose Schrift (Helvetica) wie im Muster.
\newcommand{\hsfont}{\fontfamily{phv}\selectfont}
\newcommand{\Schlagwoerter}[2]{\par\vspace{2ex}\hrule height 0.4pt\vspace{1ex}\noindent\textbf{#1:} #2\par}
\makeatletter\renewcommand*\@pnumwidth{2.6em}\makeatother % Platz für römische Seitenzahlen im Inhaltsverzeichnis
% \BHTTitelseite[Zusatzzeile]  – Zusatzzeile z.B. für den Auszug (optional)
\newcommand{\BHTTitelseite}[1][]{%
  \begin{titlepage}
    \hsfont\setlength{\parindent}{0pt}\setlength{\parskip}{0pt}%
    \setstretch{1.15}%
    \centering
    \vspace*{-2ex}
    \includegraphics[width=0.42\textwidth]{bht/BHT-Logo-horizontal}\par
    \vspace{3ex}
    {\large Berliner Hochschule für Technik\par}\vspace{1ex}
    {\large \Fakultaet\par}\vspace{1ex}
    {\large Studiengang \Studiengang\par}
    \vspace{7ex}
    {\LARGE\bfseries \Titel\par}
    \ifZweiteSprache\vspace{1.5ex}{\small \TitelEN\par}\fi
    \ifx\relax#1\relax\else\vspace{2ex}{\normalsize #1\par}\fi
    \vspace{7ex}
    {\Large\bfseries \Arbeitstyp\par}
    \vspace{2ex}
    zur Erlangung des akademischen Grades\par\vspace{1ex}
    {\large \Abschluss\par}
    \vspace{6ex}
    vorgelegt von\par\vspace{1ex}
    {\large \Autor\par}\vspace{0.5ex}
    Matrikelnummer: \Matrikel\par
    \vfill
    {\raggedright
    \begin{tabular}{@{}p{4.2cm}l@{}}
      Betreuer BHT:       & \Betreuer \\
      Betreuer extern:    & \BetreuerExtern \\
      Gutachter:          & \Gutachter \\
    \end{tabular}\par
    \vspace{4ex}
    \begin{tabular}{@{}p{4.2cm}l@{}}
      Eingereicht:    & \Abgabedatum \\
    \end{tabular}\par}

  \end{titlepage}}

% Verweis, der im Auszug (ohne Kapitel 3) auf einen Text ausweicht
\makeatletter
\newcommand{\crefoder}[2]{\@ifundefined{r@#1}{#2}{\cref{#1}}}
\makeatother
